\textbf{A diáspora judaica e o lugar} investiga a maneira pela qual a cultura e as tradições judaicas se apropriam dos espaços de trânsito, os movimentos cíclicos --- por meio da ``consagração'' de um lugar há a repetição da cosmogonia, ou seja, a vitória do lugar sagrado sobre o espaço profano deve ser repetida simbolicamente todos os anos, pois todos os anos o mundo necessita ser criado novamente desde o seu princípio --- e constantes das diásporas, e os tornam uma possibilidade de pertencimento. Um começo, um meio e um começo. No espaço talmúdico deste texto há a intenção de definir lugar, ainda que na ausência de um. É necessariamente neste ponto --- em que não se sabe ao certo o que se encontrará --- que imaginamos ter algo a dizer. Pois, como intérprete de uma longa tradição de textos e de livros, acredito que o nosso mundo --- e todas as coisas que ele contém --- está sendo criado e reinventado por nós a todo instante.

\textbf{Tamara Crespin} \lipsum[2]

\textbf{Paulina Cho} \lipsum[2]

